\section{Joint depth and severity of incidence witnesses}
\label{sec:v112-load}
The factorial estimate treats the witness thresholds only through
 their sum. It therefore forgets that there are at most two collision
indices at any depth. It also forgets how much an individual incidence
lies below its guard threshold. We retain both pieces of information.
All estimates below concern the same target-roof measure as
Theorem~\ref{thm:v111-roof-pressure}, with its original exact labels.
They use its inherited geometric inputs, not a new conditional
mixing or local-limit assumption.

Put
\begin{equation}\label{eq:v112-depth-parameters}
 r=q_0^{1/4},\qquad \alpha=-\log r=\tfrac14|\log q_0|,
 \qquad b_\varepsilon=4K_\star\varepsilon/\alpha.
\end{equation}
For a regular trajectory define its incidence-witness set and load by
\begin{equation}\label{eq:v112-load}
 \begin{split}
 J_\varepsilon(x)&=\{i:1\le i<m,\ c_i(x)<2\varepsilon r^{d_i}\},\\
 \mathscr L_\varepsilon(x)
  &=\sum_{i\in J_\varepsilon(x)}\log\frac{2\varepsilon}{c_i(x)}\\
  &=\alpha\sum_{i\in J_\varepsilon(x)}d_i
       +\sum_{i\in J_\varepsilon(x)}
                       \log\frac{2\varepsilon r^{d_i}}{c_i(x)}.
 \end{split}
\end{equation}
The empty sum is zero. Both summands in the last line are
nonnegative. In particular this is not a signed logarithmic weight,
and it introduces no clearance factor. The Borel witness set is
allowed to vary with the roof and angle.

For a nonnegative Borel weight $h$, write
\[
 D[h](t)=\frac{dF_\#(h\one_\Pi\nu_R^*)}{dt}(t)
\]
when this density is defined. The positive identities of the
preceding section give a common version at fixed $m,R,n,k$.
The next theorem also proves existence and finiteness for its weights.

\begin{theorem}[Exponential moment of the original incidence load]
\label{thm:v112-load}
For every finite collision count, every $\varepsilon>0$ and
$0<\lambda<1$,
\begin{equation}\label{eq:v112-mgf}
 \bigl\|D[e^{\lambda\mathscr L_\varepsilon}]\bigr\|_\infty
 \le C_\star
   \exp\left\{\frac{4K_\star\varepsilon}
                         {e^{\alpha(1-\lambda)}-1}\right\}.
\end{equation}
Consequently, for $u>b_\varepsilon$,
\begin{equation}\label{eq:v112-load-tail}
 \bigl\|D[\one_{\{\mathscr L_\varepsilon\ge u\}}]\bigr\|_\infty
 \le C_\star\min\left\{1,
               \exp[-u+2\sqrt{b_\varepsilon u}]\right\}.
\end{equation}
Both estimates are uniform in the radius, exact labels and target
roof. They remain valid with any further Borel restriction of the
original measure, including the persistent source through
Lemma~\ref{lem:v111-domination}.
\end{theorem}
\begin{proof}
Set $y_i=2\varepsilon/c_i$. At a witness, $y_i>r^{-d_i}$, whence
\[
 y_i^\lambda
 =y_i y_i^{-(1-\lambda)}
 \le y_i r^{(1-\lambda)d_i}.
\]
At every regular trajectory we therefore have the positive inequality
\begin{align*}
 e^{\lambda\mathscr L_\varepsilon}
 &=\prod_{i=1}^{m-1}
       \left[1+(y_i^\lambda-1)\one_{\{i\in J_\varepsilon\}}\right]\\
 &\le\prod_{i=1}^{m-1}
                 \left(1+\frac{2\varepsilon r^{(1-\lambda)d_i}}{c_i}\right).
\end{align*}
Apply the finite product insertion
\eqref{eq:v111-product-insertion}. There are at most two indices
at every positive depth, so
\[
 \sum_{i=1}^{m-1}2\varepsilon r^{(1-\lambda)d_i}
 \le4\varepsilon\sum_{d\ge1}r^{(1-\lambda)d}
 =\frac{4\varepsilon}{e^{\alpha(1-\lambda)}-1}.
\]
This proves \eqref{eq:v112-mgf} after the complete physical sum.
Every product is finite before this last geometric upper bound is
used; no independence of collisions is assumed.

Positivity gives
$\one_{\{\mathscr L_\varepsilon\ge u\}}
 \le e^{-\lambda u}e^{\lambda\mathscr L_\varepsilon}$.
Writing $v=1-\lambda$ and using $e^{\alpha v}-1\ge\alpha v$,
we obtain, for every $0<v<1$,
\[
 \bigl\|D[\one_{\{\mathscr L_\varepsilon\ge u\}}]\bigr\|_\infty
 \le C_\star\exp\{-u+vu+b_\varepsilon/v\}.
\]
When $u>b_\varepsilon$, take $v=\sqrt{b_\varepsilon/u}$.
The independent bound $\|D[1]\|_\infty\le C_\star$ follows from
the empty marked set and supplies the minimum in
\eqref{eq:v112-load-tail}. Borel restrictions decrease a positive
measure. For the persistent source the pointwise coefficient
inequality in Lemma~\ref{lem:v111-domination} supplies this decrease.
\end{proof}

\begin{corollary}[Accumulated depth and severe shallow incidences]
\label{cor:v112-load-consequences}
For $v>0$, each of the events
\[
 \left\{\alpha\sum_{i\in J_\varepsilon}d_i\ge v\right\},
 \qquad
 \left\{\sum_{i\in J_\varepsilon}
                  \log\frac{2\varepsilon r^{d_i}}{c_i}\ge v\right\}
\]
is contained in $\{\mathscr L_\varepsilon\ge v\}$ and has the
bound \eqref{eq:v112-load-tail} whenever $v>b_\varepsilon$.
An incidence witness satisfying $c_i\le2\varepsilon e^{-v}$ is
also in that event, irrespective of its depth.
\end{corollary}
\begin{proof}
Use the nonnegative decomposition in \eqref{eq:v112-load}. For the
last assertion, the single indicated summand is at least $v$.
\end{proof}
Thus a single sufficiently severe witness in an end layer can be
paid. The conclusion does not require either a deep witness or the
large witness count of Theorem~\ref{thm:v111-witness-payment}.
It does not assert a bound for an unvisited third-disk tangency.

\section{Quadratic tails for the number of incidence witnesses}
\label{sec:v112-count}
There is also an explicit finite-count improvement which retains
all threshold locations, rather than passing to their total sum.
The calculation is deterministic; it makes no independence claim
about the billiard or its symbolic itinerary.

For $0<r<1$ write
\[
 (r;r)_a=\prod_{j=1}^a(1-r^j),\quad (r;r)_0=1,
 \quad (r;r)_\infty=\prod_{j=1}^\infty(1-r^j)>0,
 \quad f(s)=\left\lfloor\frac{(s+1)^2}{4}\right\rfloor.
\]
The infinite product is positive because the tail of
$\sum_j-\log(1-r^j)$ is bounded by a constant times $\sum_jr^j$.

\begin{lemma}[The two-end depth polynomial]\label{lem:v112-depth-poly}
For $0\le s\le m-1$, put
$E_{m,s}(r)=e_s((r^{d_i})_{1\le i<m})$. Then
\begin{equation}\label{eq:v112-poly-bounds}
 r^{f(s)}\le E_{m,s}(r)
 \le\sum_{a=0}^s
    \frac{r^{a(a+1)/2+(s-a)(s-a+1)/2}}
         {(r;r)_a(r;r)_{s-a}}
 \le\frac{(s+1)r^{f(s)}}{(r;r)_\infty^2}.
\end{equation}
For $s>m-1$ the finite polynomial coefficient is zero and only the
upper inequalities are asserted. In particular the coefficient has
quadratic logarithmic rate, uniformly over $m\ge s+1$:
\begin{equation}\label{eq:v112-poly-sharp}
 \log E_{m,s}(r)=-\frac{\alpha}{4}s^2+O_r(s+\log(s+1)).
\end{equation}
This is sharpness of the depth polynomial, not a lower bound for a
physical roof density.
\end{lemma}
\begin{proof}
The multiset of finite collision depths is contained in two copies
of $1,2,3,\ldots$. Coefficientwise positivity bounds its generating
polynomial by $\prod_{d\ge1}(1+zr^d)^2$.
To compute one factor, for $a\ge1$ write
$1\le d_1<\cdots<d_a$ as $d_j=j+h_j$, where
$0\le h_1\le\cdots\le h_a$. Setting
$g_1=h_1$ and $g_j=h_j-h_{j-1}$ gives
\[
 \sum_{1\le d_1<\cdots<d_a}r^{d_1+\cdots+d_a}
 =r^{a(a+1)/2}
    \prod_{j=1}^a\sum_{g_j\ge0}r^{(a-j+1)g_j}
 =\frac{r^{a(a+1)/2}}{(r;r)_a}.
\]
The empty factor contributes one. Multiplying the two positive
series proves the middle sum in \eqref{eq:v112-poly-bounds}; all
infinite passages are monotone limits of nonnegative coefficients.
The minimum of its exponent over $a$ is $f(s)$, attained by the
two most balanced choices. Each denominator is at least
$(r;r)_\infty^2$, giving the last inequality.

For the lower bound take the $s$ smallest available depths. If
$s=2k$ these consist of two copies of $1,\ldots,k$ and have sum
$k(k+1)$. If $s=2k+1$ take these and one copy of $k+1$, with sum
$(k+1)^2$. These depths are available whenever $m\ge s+1$, including
the single central depth when $m$ is even. Their product is one
term of $E_{m,s}$. Taking logarithms of the two bounds proves
\eqref{eq:v112-poly-sharp}.
\end{proof}

\begin{theorem}[Quadratic witness-count tail at the target roof]
\label{thm:v112-count}
With $N^\varepsilon=|J_\varepsilon|$, for every $s\ge1$,
\begin{equation}\label{eq:v112-count-tail}
 \bigl\|D[\one_{\{N^\varepsilon\ge s\}}]\bigr\|_\infty
 \le\frac{C_\star(s+1)(2K_\star\varepsilon)^s}
               {(r;r)_\infty^2}\,r^{f(s)}.
\end{equation}
The event is empty for $s>m-1$. The same estimate holds for the
persistent source. For every fixed $B$ and every $\eta>0$, take
\begin{equation}\label{eq:v112-root-log-count}
 \sigma_m=\left\lceil
       2\sqrt{\frac{(2+\eta)\log(m+1)}{\alpha}}\right\rceil.
\end{equation}
Then the source with at least $\sigma_m$ incidence witnesses has
zero $m^2$-normalized essential height, with finite-count rate
$m^{-\eta+o_B(1)}$.
\end{theorem}
\begin{proof}
Use the first, unrelaxed inequality in
\eqref{eq:v111-factorial}, with thresholds
$2\varepsilon r^{d_i}$, and apply Lemma~\ref{lem:v112-depth-poly}.
Since $f(s)\ge s^2/4$, the logarithm of $m^2$ times the upper bound
at $s=\sigma_m$ is at most
\[
 -\eta\log(m+1)
 +\sigma_m\log(2K_\star\varepsilon(B))
 +\log(\sigma_m+1)
 +\log\frac{C_\star}{(r;r)_\infty^2}.
\]
For fixed $B$ the last three terms are $O_B(\sqrt{\log(m+1)}
+\log\log(m+2))$. This proves the claim. No outer width is changed.
\end{proof}
The square-root logarithmic threshold is strictly smaller in order
than the logarithmic-over-logarithmic threshold of the earlier
factorial bound. The estimate depends on the actual two-end depth
geometry, not on additional cancellation or a count of physical words.
